\documentclass[english,a4paper]{my_memo}
\usepackage{booktabs,tabularx}
\setlength{\emergencystretch}{3em}
\hypersetup{colorlinks=true,urlcolor=blue,linkcolor=blue}
\newcommand{\code}[1]{\nolinkurl{#1}}
\title{Accelerating Physical Marching in EXHALE v1.10}
\author{Implementation and literature assessment}
\date{November 9, 2026}
\begin{document}
\maketitle
\begin{abstract}
Implicit hydrodynamics is a plausible way to accelerate acoustic-CFL-limited EXHALE marching, but neither unconditional stability nor a larger timestep guarantees a faster accurate solution. Dorfi (1999) provides an early conceptual basis; Leidi (2022) and Viallet (2013, 2016) supply more direct evidence on implicit compressible integration and its cost. Ghosh and Menon (2010) explicitly retain a Courant timestep restriction in their implementation, so that paper does not demonstrate the particular acceleration sought here. The preferred development candidate is fixed-grid conservative implicit compressible hydrodynamics, comparing direct radial linear algebra with acoustic preconditioning and retaining consistent species and thermal coupling. First measure the existing phase costs and the cells controlling the timestep. No speedup factor has been measured in this study, and the frozen campaign must not be changed on this evidence alone.
\end{abstract}

\section{Scope and evidence}
This report evaluates methods for reducing the wall time required to march EXHALE v1.10 over a specified physical interval while preserving the required accuracy. It reviews thirteen final journal articles and checks their applicability against the solver source. Published timings, arithmetic ratios derived from those timings, and proposed EXHALE methods are identified separately. No EXHALE performance benchmark was run for this report.

Implementation statements refer to the inspected working-tree source, whose hashes are recorded in the accompanying evidence directory. They do not establish the configuration of a frozen campaign executable. Compilation time and the presence of build products are not performance criteria for this investigation.

\subsection{Quantitative evidence at a glance}
\begin{center}
\begin{tabularx}{\textwidth}{@{}p{0.20\textwidth}p{0.21\textwidth}X@{}}
\toprule
Study & Published result or derived ratio & What was compared\\
\midrule
Leidi (2022) & 10--20 times faster & Implicit versus explicit time integration of a periodic vortex at $\mathcal M\sim10^{-3}$, on the same grid.\\
Viallet (2013) & About 2--3 and 15--20 times faster & Stellar models at characteristic Mach numbers about 0.1 and 0.01; explicit MUSIC comparators are not optimized.\\
Viallet (2016) & Progress ratios 0.51, 5.77, 105.3 & Physical time reached in six wall-clock hours at initial Mach numbers $10^{-1}$, $10^{-2}$, $10^{-3}$; different final physical times.\\
Knoll (2001) & CPU ratios 7.25 and 17 & Nonlinear second-order versus lagged semi-implicit radiation diffusion at comparable error, without material motion.\\
Schlegel (2012) & Derived speed ratios 1.41--2.44 & Multirate versus reference integration in different advection/reaction tests; not a range for one wind problem.\\
Esent\"urk (2018) & Whole-model time reduced 1.81--3.59\% & Selected quasi-Newton chemistry corrections; chemistry itself saves 12.74--13.48\%.\\
\bottomrule
\end{tabularx}
\end{center}
These results do not determine an EXHALE speedup range. Dorfi's large CFL numbers measure timestep size, and Townsend's timings measure scalar cooling updates; neither is a complete wind-runtime comparison. The detailed sections below give the methods, assumptions, and limitations behind the numbers.

\section{What currently limits marching}
\subsection{Hydrodynamic timestep and update}
\code{src/modules/time_step/eval_dt.f90}, routine \code{eval_dt}, computes
\begin{equation}
 \Delta t_j=C_{\rm CFL}\frac{V_j}{\max(A_{j-1/2}s_{j-1/2},A_{j+1/2}s_{j+1/2})},
 \qquad \Delta t=\min_j\Delta t_j,
\end{equation}
where $A=r^2$, $V=(r_{j+1/2}^3-r_{j-1/2}^3)/3$, and the face signal is the maximum neighboring $|v|+c_s$. The caloric EOS supplies the adiabatic index when active. Thus fine cells with large sound speed can determine the global step even when the material velocity is small. This is a verified algorithmic restriction; the controlling cell and actual Mach number have not been measured for a campaign state here.

The main program calls this routine near line 2081 and performs explicit three-stage SSP-RK hydrodynamics starting near line 2304. Species advection uses the corresponding face mass fluxes, rather than an independent velocity estimate. Positivity checks, flux corrections, and rejected-step reductions constrain the usable interval further. Accepted physical time advances by $\Delta t R_0/v_0$ near line 3488. Increasing an input CFL or prescribing a larger diagnostic timestep does not make these explicit stages implicit.

Local time intervals already exist in \code{eval_dt}, but they describe pseudo-time relaxation when different cells use different intervals. Such intervals cannot be interpreted as a common physical clock for the final three flow-crossing times. They therefore cannot directly replace the campaign marching trajectory.

\subsection{Thermal and compositional coupling}
Near lines 2708 and 2779 the main program invokes \code{ioniz_eq} and, conditionally on \code{use_semi_implicit_energy}, \code{solve_energy_semi_implicit}. The latter solves a local thermal-energy equation using a bracketed Newton--bisection iteration. It retains the old thermal energy and closes the EOS at the supplied composition and heating. It is not an implicit hydrodynamic pressure solve, nor a fully simultaneous global radiation--chemistry--hydrodynamics solve. The explicit energy option also remains callable.

The radiation implementation in
\code{src/modules/radiation/ionization_equilibrium.f90} includes column attenuation and state-dependent local optical depths. Changing a cell can therefore affect radiation elsewhere. A fully coupled Jacobian is not automatically limited to the nearest neighboring hydrodynamic cells. Chemistry, opacity, and the caloric EOS must remain consistent with the chosen implicit state.

The source in \code{src/modules/time_step/steady_newton.f90} contains nonlinear-solver routines and banded LAPACK interfaces that may be useful in a physical-time integrator. Reuse requires a new time-dependent residual and checks of its boundary derivatives and source coupling; the existing stationary residual is not a physical-time update.

The invoked \code{ioniz_eq} path solves composition through equilibrium relations coupled to radiation and the configured transport constraints. It is not simply the time-dependent reaction-network ODE used in the atmospheric chemistry papers. Reducing repeated nonlinear work while converging the same closure is a numerical optimization. Replacing that closure by a kinetic ODE changes the physical model, even if a published ODE solver is faster. This distinction must be preserved when assessing chemical subcycling or ASIS.

\section{Implicit hydrodynamics and radiation coupling}
\subsection{Dorfi (1999)}
Dorfi\ \cite{dorfi} evolves spherical radiation hydrodynamics on an adaptive moving grid using an implicit nonlinear system. The discussion around Eqs.~(41)--(46), pp.~166--167, explicitly contrasts the acoustic CFL restriction with implicit integration. The Newton system has a block pentadiagonal structure for that paper's stencil; the reported factorization cost scales as $M^3N$ for $M$ variables at each of $N$ grid locations. Timestep selection is still restricted by nonlinear convergence and permitted state changes.

A quantitative example appears on p.~165: the shock-tube evolution over $0\leq t\leq1$ takes about 100 steps, with typical Courant numbers $10^2$--$10^3$. These are published values, not measurements repeated here. They imply steps roughly 100--1000 times the acoustic CFL interval as defined in that example, not a measured 100--1000-fold reduction in wall time. The moving adaptive grid follows the wave structure, so a fixed-grid EXHALE calculation cannot inherit this ratio without validation. No matched explicit-versus-implicit wall-time benchmark is established by this example.

This supports removing an explicit acoustic stability restriction, not taking arbitrary accurate steps. The discussion around Eq.~(21) also emphasizes temporal damping and centering for oscillations. EXHALE should borrow the conservative implicit-time principle, while treating the adaptive mesh and the radiation moment formulation as separate design choices. A moving mesh would require changing fluxes, geometric conservation, boundary handling, species transport, restarts, and comparison sampling. It is not necessary for a first acceleration experiment.

\subsection{Ghosh and Menon (2010)}
Ghosh and Menon\ \cite{ghosh} use Lagrangian meshes, a tridiagonal velocity solve, and nested pressure and temperature iterations coupled to discrete-ordinates radiation transport. Their algorithm description on p.~7495 explicitly selects the timestep using a Courant condition and an additional compressibility restriction. The convergence study reports approximately first-order spatial and temporal behavior for the tested discretization.

The relevant contribution is a concrete implicit coupling and verification procedure. It is not evidence that retaining their timestep controller will accelerate EXHALE beyond its acoustic CFL interval. Their radiation-driven shocks and point explosions do not reproduce a photoionized escaping atmosphere. In particular, a moving material mesh in an open wind needs treatment of mass entering and leaving the domain. Replacing EXHALE's radiation model with theirs solely for speed has no demonstrated physical or computational benefit.

\subsection{Independent IMEX evidence}
Gardner et al.\ \cite{gardner} compare implicit--explicit partitions, time integrators, and nonlinear solvers for nonhydrostatic atmospheric equations. Their tests show that allowing larger stable steps can fail to repay the cost of a more extensive implicit solve. Their gravity-wave and baroclinic tests are not EXHALE benchmarks. They support measuring total cost at matched error and choosing the implicit terms carefully; no numerical speedup from that paper is transferred here.

\section{Additional acceleration strategies}
\subsection{Conservative multirate physical integration}
Schlegel et al.\ \cite{schlegel} combine explicit advection with implicit diffusion/reaction and different physical timesteps in different spatial regions. Their methods section discusses conservation and the nested integration construction. This differs from assigning an independent pseudo-time to each cell: all regions must reach common physical output times with consistent integrated exchanges across interfaces.

Their final journal article, p.~1404, reports a 59\% average cost reduction for pure advection in the academic test, about 36\% for its full advection--diffusion--reaction system, and 29\% average reduction in the realistic case (Fig.~9). Converting a fractional reduction $d$ into a speed ratio gives $1/(1-d)$: approximately 2.44, 1.56, and 1.41, respectively. These are conversions of published timings, not EXHALE predictions. The paper attributes part of the coupled-system benefit to fewer changes in the implicit source forcing and fewer small steps or restarts. It does not demonstrate that infrequent chemistry updates are accurate in a photoionized planetary wind.

For EXHALE, a spatial multirate method is attractive only if a small fraction of cells determines the global interval. Accumulate the same time-integrated mass, momentum, energy, and species flux on both sides of every interface. Radiation column changes remain nonlocal and may require more frequent synchronization than material transport. Existing \code{dt_loc} values alone do not implement this scheme. A simple work estimate illustrates the opportunity: if a fraction $f$ of cells needs the fine step and all other cells can use $m$ times that step, the ideal work ratio relative to uniform fine stepping is
\begin{equation}
 S_{\rm ideal}=\frac{1}{f+(1-f)/m}.
\end{equation}
For the hypothetical $f=0.1$, $m=10$, this is 5.26. Interface work, radiation synchronization, implicit iterations, and parallel imbalance reduce that gain. The global clock and all-cell K=3 sampling must be preserved.

\subsection{Acoustic subcycling with less frequent expensive sources}
A related alternative is to retain small explicit hydrodynamic steps inside a larger thermal/radiation interval. This does not remove acoustic work, but can save repeated expensive source calculations when source evolution is sufficiently slow. If source work occupies a fraction $f_Q$ of the original cost and is evaluated once for every $m$ hydro steps, the optimistic speed ratio is
\begin{equation}
 S_{\rm ideal}=\big[(1-f_Q)+f_Q/m\big]^{-1}.
\end{equation}
This estimate is conditional on unchanged hydro cost and negligible coupling overhead. Use interpolation or a controlled coupled approximation rather than holding rapidly varying photoheating fixed. A full interval versus subdivided intervals must verify thermal and species errors. This is an engineering candidate motivated by the multirate and IMEX literature, not an implemented EXHALE option or a published EXHALE performance result.

\subsection{Linearly implicit stages and Jacobian reuse}
Gardner et al.\ \cite{gardner} examine a Rosenbrock-like linearized solve as well as Newton iteration. Linearizing an implicit stage can lower its cost, but may require smaller steps to preserve accuracy. For EXHALE, reusing a Jacobian or factorization over sufficiently similar stages could reduce repeated differentiation and linear algebra. Refresh it after rejected steps, substantial opacity/EOS changes, or slow residual reduction. Do not reuse a stale factorization as though the nonlinear equations had been solved. This is a way to make the proposed physical-time implicit scheme cheaper, not a new stationary solution method.

\subsection{Low-Mach fluxes and hydrostatic balance}
Reducing artificial dissipation and maintaining pressure--gravity balance can improve accuracy of a slow flow. Neither change alone removes the explicit acoustic stability interval. The current \code{Num_Fluxes.f90} already contains a selectable \code{low_mach_velocity_jump} branch, read by \code{input_read.f90}, and well-balanced branches. Its existence must not be counted as a new acceleration implementation; the numerical effect and active campaign settings need separate verification.

A second option, \code{low_mach_dissipation.f90}, adds a fourth-difference stress. It is not acoustic preconditioning or physical viscosity. Its variable-coefficient stress is explicitly documented as sign-indefinite, so simply increasing this damping to obtain a settled profile is not a physically justified acceleration. 

Pure low-Mach or anelastic equations filter acoustic dynamics by changing the governing approximation. They are not suitable as a global replacement when the wind crosses the sonic point. A low-Mach inner region coupled to a compressible outer region would require a separate interface and validity study. Prefer an all-Mach compressible method for the first investigation.

\subsection{Diffusion acceleration is not the current primary opportunity}
Super-time-stepping extends explicit stability along the negative real axis for suitable diffusion operators; it does not generally stabilize acoustic eigenvalues. More decisively, the inspected EXHALE path already calls \code{viscous_conduction_step} near main-program line 3078. That routine assembles and solves tridiagonal momentum and thermal systems via \code{solve_tridiagonal} (calls near lines 1110 and 1177). Its linear solve costs scale linearly with radial cell count. Replacing this existing implicit diffusion treatment by multiple explicit stages has no demonstrated advantage. Consider that replacement only after measuring this particular phase and preserving positivity, thermal accuracy, and boundary energy exchange.

\section{Gravity, low-Mach flow, and nonlinear source coupling}
\subsection{Edelmann et al. (2021): preserve the small physical imbalance}
Edelmann et al.\ \cite{edelmann} compare three well-balanced formulations for fully compressible Euler flow with gravity: Cargo--LeRoux, $\alpha$--$\beta$, and deviation well-balancing (Sect.~4). The latter two preserve small perturbations more accurately in strongly stratified, low-Mach tests. Their implicit ESDIRK23 integration and low-dissipation flux are distinct ingredients: balancing gravity alone does not remove the acoustic timestep restriction. The one-dimensional hydrostatic tests retain extremely small velocities, with Mach numbers below about $10^{-12}$ in the reported tests; this is an equilibrium-preservation result, not a CPU acceleration factor. The two-dimensional tests also expose pressure--velocity decoupling and difficulties with unresolved internal gravity waves. In the Keplerian disk test (Sect.~5.4), with Mach numbers around 0.6--1, explicit integration is more efficient. Thus the paper supports careful spatial balance, not a universal advantage for implicit integration.

Deviation well-balancing uses a specified stationary reference. Subtracting the residual of an arbitrary marching snapshot would generally change the equations and could artificially pin the flow. A hydrostatic reference does not cancel real photoheating, reactions, or their subsequent pressure response. EXHALE should preserve the pressure--gravity balance while retaining these physical sources; no additional stationary solution is required. The long equilibrium tests in the paper do not replace the campaign's K=3 acceptance window. 

\subsection{Leidi et al. (2022): implicit compressible integration}
Leidi et al.\ \cite{leidi} combine a low-dissipation LHLLD flux, deviation well-balancing, and implicit--explicit Strang splitting (IESS). The partition in Sect.~4, Eq.~(48), is important: explicit induction half-steps surround a \emph{fully implicit density, momentum, and energy update}, using ESDIRK23. It is not simply a pressure-only IMEX update. In a nonmagnetic application the induction substeps disappear, leaving implicit compressible hydrodynamics. The timestep uses fluid and Alfv\'en velocities rather than the fast magnetosonic speed. Their wave tests explicitly show that stable large steps need not accurately resolve fast waves.

The Balsara-vortex timing experiment in Sect.~5.2 and Fig.~8 uses a $40\times40$ periodic grid, the same spatial discretization for both integrators, and no snapshot output. The authors reduce the explicit stability step by 20\%. They report that IESS becomes faster near initial maximum sonic Mach numbers $\mathcal M\simeq0.04$ for $\beta_K=0.1,1$, and $\mathcal M\simeq0.02$ for $\beta_K=10$; at $\mathcal M\simeq10^{-3}$ it is about 10--20 times faster. Here $\beta_K$ is a magnetic-to-rotational-kinetic-energy ratio, not plasma beta. These are published same-grid timings with accompanying accuracy tests, not a matched-error EXHALE benchmark. The periodic ideal-gas vortex has neither planetary gravity nor the wind's radiation and chemistry.

The Kelvin--Helmholtz tests show a separate benefit from lower spatial dissipation: less resolution may be needed. A source inconsistency must be retained in the assessment: Sect.~5.3 assigns the twofold and fourfold resolution increases for HLLD relative to LHLLD to Mach numbers $10^{-2}$ and $10^{-3}$, respectively, whereas Sect.~6 reverses those associations. We do not use the reversed summary or transfer the paper's two-dimensional resolution-cost estimates to a one-dimensional wind. 

\subsection{Bispen et al. (2017): useful pressure elimination, restricted physics}
Bispen et al.\ \cite{bispen} split a linear acoustic--gravity operator from nonlinear transport and use BDF2 or ARS(2,2,2) integration. Eliminating momentum and density gives a pressure-related elliptic equation; their implementation uses UMFPACK. This suggests an inexpensive radial pressure solve, but their derivation assumes a constant-gravity ideal gas and a potential-temperature formulation for smooth adiabatic flow. In Sect.~3, Eqs.~(15)--(19), they neglect a nonlinear pressure contribution of order $\epsilon^4$ in the adopted low-Mach scaling. This is not a justified omission in a heated, composition-changing, transonic atmosphere.

The asymptotic analysis in Sect.~4 specializes to constant background potential temperature and periodic conditions. The stability analysis in Sect.~5 explicitly omits gravity. Neither proves stability for EXHALE's spherical gravity, open boundaries, caloric EOS, and reacting gas. Large reported CFL numbers establish the possibility of large stable steps; they are not wall-time speed ratios. Their acoustic-pulse and warm-bubble tests also show that resolving the fast dynamics can require smaller steps. The useful transferable element is the elimination strategy, after deriving the correct energy and EOS coupling, or its use as an approximate preconditioner for a full residual.

\subsection{Knoll et al. (2001): nonlinear closure and temporal order matter}
Knoll et al.\ \cite{knoll} study one-dimensional grey radiation diffusion coupled to material energy. Material motion and material heat conduction are excluded (Sect.~2). Their second material model includes ionization energy through a Saha-type treatment, but that ionization does not alter the opacity. Their semi-implicit method (SI) lags nonlinear coefficients; NK1 solves backward Euler nonlinearly, and NK2 solves Crank--Nicolson nonlinearly. Newton--Krylov is used within \emph{physical time steps}, not to require a stationary solution.

For comparable error in Problem~1 (p.~32), the normalized CPU costs are SI: 7.25, NK1: 1.7, NK2: 1. For Problem~2 (p.~33), they are 17, 4, and 1. Thus NK2 gives published cost reductions corresponding to speed ratios of 7.25 and 17 relative to SI. The corresponding thermal-front CFL ratios of approximately 200 and 400 are not CPU gains and are not acoustic CFL numbers. In Problem~3, Table~2 gives SI cost 25.9 with error 0.59, versus NK2 cost 1 with error 0.044. This is \emph{not} an equal-error 25.9-fold comparison. It demonstrates why an apparently unchanged result under timestep reduction can be misleading when nonlinear coefficients are insufficiently converged.

For EXHALE, this motivates convergence of radiation, composition, and thermal energy at the accepted state. It does not establish that EXHALE currently performs only the paper's lagged SI update: the inspected source already has a coupled iteration. Nor does the paper prove that Crank--Nicolson is always positive or nonoscillatory; its implicit-midpoint variant encountered convergence problems. Temporal order and damping must be verified for the complete split update.

\section{Reducing implicit-solver and source costs}
\subsection{Viallet et al. (2013): spend less on each implicit step}
Viallet et al.\ \cite{viallet13} compare nonlinear solvers for Crank--Nicolson compressible hydrodynamics, including gravity and radiative diffusion. The central equation is
\begin{equation}
 F(U^{n+1})=\frac{U^{n+1}-U^n}{\Delta t}
 -\frac{R(U^{n+1})+R(U^n)}{2}=0.
\end{equation}
The time term is essential: this is marching, not $R(U)=0$. Their best two-dimensional results use Broyden updates and an incomplete-LU preconditioner built at the first nonlinear iteration and reused within that time step. A relatively loose linear solve can save work while a sufficiently strict nonlinear convergence test preserves accuracy. The paper's tolerances cannot be copied to differently scaled EXHALE residuals or trace abundances.

Section~5.2 reports roughly 2--3-fold gains for a red-giant model and 15--20-fold gains for a young-Sun model relative to the explicit methods considered. Their characteristic convective Mach numbers are about 0.1 and 0.01. The paper explicitly cautions that its explicit MUSIC implementations are not optimized competitors: Adams--Bashforth is restricted to a particularly small stable CFL. The vortex study finds temporal error becomes important as material transport crosses multiple cells in a step. Its $\mathrm{CFL}_{\rm adv}\lesssim2$ accuracy observation is specific to that idealized test, not an EXHALE limit.

The practical lesson is to compare Jacobian assembly, factorization, residual evaluation, and iteration costs separately. Broyden can win despite needing more nonlinear iterations. For a short one-dimensional radial mesh, a directly factored block-banded hydro system may be cheaper than a matrix-free iteration, unlike the large multidimensional systems motivating MUSIC. This is an EXHALE design inference, not a timing result in the paper. Nonlocal radiation coupling must still be converged outside the banded hydro block or represented in the full system.

\subsection{Viallet et al. (2016): approximate the solve, preserve the equations}
Viallet et al.\ \cite{viallet16} construct semi-implicit acoustic and acoustic--diffusion solves and use them as right preconditioners for the full nonlinear Crank--Nicolson residual (Sects.~3--4). The acoustic solve eliminates velocity into a pressure equation of Helmholtz type. Schematically, for a uniform background,
\begin{equation}
 \left[I-c_s^2\Delta t^2\nabla^2\right]\delta p=b.
\end{equation}
For nonuniform spherical flow its coefficients, geometric operators, gravity, and boundary conditions must be rederived. The acoustic--diffusion version couples pressure and internal-energy corrections. Their primitive/conserved-variable transformations and EOS derivatives are explicit parts of the algorithm (Sect.~3.2 and Appendices A--B).

The simplified solve by itself is only first order and comparatively inaccurate. Used \emph{as a preconditioner}, it changes the rate of convergence, while the converged full residual determines the physical update. This distinction provides a useful way to use an approximate low-Mach pressure operator without dropping nonlinear pressure or thermochemical terms from EXHALE's governing equations. It is not completely matrix-free: the smaller preconditioner matrix is formed and stored. Omitting stiff source derivatives from it may make it ineffective, even when the full residual remains correct.

Table~1 compares physical time reached in a fixed six-hour wall-clock allocation for a three-dimensional Taylor--Green vortex. The explicit Adams--Bashforth calculation uses $\mathrm{CFL}_{\rm hydro}=0.1$. The following ratios are calculated here from the published table, not from reruns:
\begin{center}
\begin{tabular}{@{}rrrr@{}}
\toprule
Initial Mach & Implicit final time & Explicit final time & Ratio\\
\midrule
$10^{-1}$ & 17.8 & 34.8 & 0.51\\
$10^{-2}$ & 20.9 & 3.62 & 5.77\\
$10^{-3}$ & 31.6 & 0.300 & 105.3\\
$10^{-4}$ & 55.1 & 0.0365 & 1509.6\\
\bottomrule
\end{tabular}
\end{center}
These are progress ratios at equal wall time, with different final physical times, not costs of reaching the same final solution at matched error. They depend on a periodic ideal-gas vortex, the explicit comparator, and solver settings. They cannot be used to predict a 100--1500-fold EXHALE gain. At $\mathcal M=10^{-6}$, reducing the implicit advective CFL from approximately 0.49 to 0.05 increases the reported physical time reached from 0.6 to 4.5, because the mean Krylov iterations decrease from 291.5 to 7.0. Even for an implicit method, \emph{a smaller step can advance physical time faster}.

The stellar comparisons are equally important. In the two-dimensional young-Sun test at advective CFL 0.5, Table~4 gives simulated-time/wall-time ratios of 74 for Broyden+ILU and 39 for the best listed JFNK+PBP variant. The older method is about 1.90 times faster. The three-dimensional red-giant comparison also changes ranking with timestep. Thus physics-based preconditioning is a candidate, not a reason to impose JFNK on EXHALE. The method's advantage in reducing multidimensional storage is weaker for a one-dimensional banded problem. The paper does not demonstrate a substantial thermal-diffusion-preconditioning gain in its stellar tests because their diffusion stiffness is modest.

\subsection{Townsend (2009): exact scalar cooling, not exact thermochemistry}
Townsend\ \cite{townsend} integrates the isochoric, optically thin scalar equation $dT/dt=-A\Lambda(T)$ exactly through a temporal evolution function and its inverse (Sect.~3.3 and Appendix~A). Density, composition, and molecular weights are fixed during that cooling substep; the adopted cooling depends on temperature alone. The paper explicitly leaves extension to evolving ionic abundances unresolved. The tests use a monatomic gas and a collisional-equilibrium cooling curve, rather than an irradiated molecular wind.

Table~1 and Sect.~4 give the cost of a cooling update on prescribed postshock states, not of a complete hydrodynamic evolution. Exact integration costs about twice the first-order explicit update, while the tested implicit methods cost about 2.5--12 times that explicit update. Exactness refers to the scalar semidiscrete equation and the chosen cooling curve. It does not remove hydrodynamic splitting error, opacity lag, or uncertainty in the cooling model.

A particularly relevant warning is in Sect.~3.2 and Fig.~3: the backward-Euler cooling equation has five roots in one illustrated case. A small nonlinear residual alone cannot identify the physically continuous thermal trajectory. This does not establish multiple roots in EXHALE's current cooling function. It motivates checking timestep refinement and branch continuity when extending its thermal timestep.

For fixed density and fixed composition $x$, a possible generalization can be derived directly from EXHALE's thermal equation:
\begin{equation}
 C_V(T;x)\frac{dT}{dt}=H(T;x)-C(T;x),\qquad
 \Delta t=\int_{T_n}^{T_{n+1}}
 \frac{C_V(T;x)}{H(T;x)-C(T;x)}\,dT.
\end{equation}
Here $C_V$ is heat capacity per unit volume. This is our mathematical observation, not a method implemented or benchmarked by Townsend. A tabulated monotone integral could avoid repeated root iterations only within a fixed source state and a connected thermal branch. A zero of $H-C$ requires special equilibrium treatment; the integral must not jump across it. Changing composition, density, radiation, or molecular excitation changes the integral. The table-generation and interpolation cost may exceed savings, so this remains lower priority than profiling the current source solve. Replacing the caloric EOS by constant $\gamma$ to obtain a convenient analytic formula would be physically unjustified.

\subsection{Esent\"urk et al. (2018): reuse chemistry linear algebra carefully}
Esent\"urk et al.\ \cite{esenturk}, Sects.~2.3--2.4, add selected quasi-Newton corrections to a backward-Euler chemistry solver, reusing the factored matrix rather than reconstructing the Jacobian for every correction. Their preferred configuration adds corrections after the second and third Newton iterations; applying them more frequently is not always better. They retain convergence checks and step reductions.

Table~2 reports that one quasi-Newton correction costs 27.3\% of a full Newton iteration in their box model. This is not a 3.66-fold speedup of the complete chemistry calculation. Table~5 reports chemistry time reductions of 12.74--13.48\% and whole-model reductions of 1.81--3.59\% in the tested configurations. The latter correspond to speed ratios of approximately 1.018--1.037. These published results distinguish savings inside a solver from savings in a complete simulation. The authors' concentration comparison with identical dynamics disables radiative feedback from the chemical species; that condition differs substantially from EXHALE's thermal coupling.

EXHALE already uses analytic-Jacobian Newton solves for some ionization systems and finite-difference derivatives for others, and its line search already tests corrections using the current Jacobian. Therefore reuse is not wholly absent. Further reuse across accepted nonlinear iterations is a specific experiment, conditional on measured Jacobian/factorization cost. Preserve the existing tests for each unknown, including trace species, and verify the final residual using the actual state. Do not replace these tests with UKCA's looser numerical tolerance or copy its dimensional secant formula without a scaling and small-denominator audit.

\subsection{Cariolle et al. (2017): adaptive chemical substeps}
The ASIS method of Cariolle et al.\ \cite{cariolle} uses a reaction-dependent semi-implicit linearization and adaptive chemical substeps (Sects.~2.2--2.3). The reaction weights are constructed consistently across species, supporting conservation, and an inexpensive concentration-curvature indicator chooses the substep before the coupled linear solve. The method is first order. Their Fortran FLUX experiment reports approximately half the computation time when a preconditioned GMRES solve replaces DGESV (Sect.~3.1); this is a comparison \emph{within ASIS}, not a factor-two gain over every chemistry integrator.

The Earth and Mars examples show that day--night transitions and departures from equilibrium need short chemical steps, while smooth periods allow longer ones. This supports measuring the distribution of source difficulty, rather than assigning one arbitrarily reduced chemistry frequency to all cells. It does not justify importing the paper's relative tolerances of 0.01--0.025 into EXHALE. Its chemical evolution equation and fixed external meteorological/photolysis updates differ from EXHALE's coupled equilibrium/transport closure. Adopting ASIS as a replacement would require a separate physical-model decision; the inexpensive error-control idea can be studied without changing the reaction model.

\section{Application to EXHALE}
The literature supports testing conservative implicit compressible hydrodynamics, with either direct block-banded linear algebra or an acoustic preconditioner, before committing to a pressure-only approximation. Leidi supplies a relevant fully compressible time integrator; Viallet supplies the distinction between the full residual and its inexpensive approximate solve. Bispen supplies useful elimination algebra but a less transferable physical approximation. These are different constructions, not interchangeable versions of one method.

\subsection{Existing discretization to preserve}
In \code{RK_rhs.f90}, lines 240--283, pressure--gravity balance is already represented through equilibrium pressure departures. A zero momentum source in that branch follows algebraic cancellation; it does not imply missing gravity. The energy flux also includes
\begin{equation}
 D_{\Phi,j}=\frac{A_+F_{\rho,+}(\Phi_+-\Phi_j)
                   -A_-F_{\rho,-}(\Phi_--\Phi_j)}{V_j}.
\end{equation}
Together with the mass update, this makes the fixed-potential update of $E_j+\rho_j\Phi_j$ use flux $F_E+\Phi_f F_\rho$. Thus the gravitational-energy consistency emphasized by Edelmann is already present in this discrete form. Adding another gravitational-work term would double count it. This algebraic observation is not a numerical conservation measurement.

In \code{Num_Fluxes.f90}, lines 24--38 and near line 394, \code{low_mach_velocity_jump} modifies only Roe acoustic dissipation. It does not modify HLLC or remove acoustic signal speeds. Leidi's LHLLD accuracy gains therefore cannot be assigned to the current HLLC calculation merely by enabling this option.

The main program's temperature--composition iteration invokes \code{ioniz_eq} and \code{solve_energy_semi_implicit}, and tests paired source-state errors near lines 2865--2873, with \code{csm_T_tol} and \code{csm_comp_tol} both $10^{-8}$. The local thermal solver retains the old energy and uses the caloric heat capacity. Knoll's results support converging this coupling, not declaring it missing. Diffusion already uses an implicit tridiagonal solve. No evidence here supports replacing it with super-time-stepping.

\subsection{Conditions that prevent direct transfer of published gains}
\begin{center}
\begin{tabularx}{\textwidth}{@{}p{0.23\textwidth}X@{}}
\toprule
Evidence & Main mismatch with EXHALE marching\\
\midrule
Leidi vortex & Periodic 2D ideal-gas MHD; no wind boundaries, chemistry, or photoheating.\\
Viallet vortices & Cartesian periodic ideal gas; explicit Adams--Bashforth comparator; no reactive radiation.\\
Viallet stellar models & Multidimensional convection, stellar EOS and radiative diffusion; surface treatment differs from an escaping atmosphere.\\
Knoll fronts & Grey radiation diffusion without material motion; thermal-front CFL, not acoustic CFL.\\
Townsend cooling & Fixed-composition scalar cooling update; no simultaneous photoheating/composition evolution.\\
Esent\"urk / Cariolle & Terrestrial or Martian chemical networks and forcing; different source closure, error tolerances, and parallel workload.\\
\bottomrule
\end{tabularx}
\end{center}
No unknown input condition is assumed to match. Published ratios are evidence that acceleration is possible, not a confidence interval for EXHALE. They must not be multiplied together: the methods can remove overlapping costs and change each other's bottlenecks.

\subsection{Accuracy constraints on a larger step}
A full implicit compressible update and an acoustic IMEX split have different remaining restrictions. An IMEX method with explicit material transport retains a material-advection stability limit. A fully implicit hydro update need not retain that same stability limit, but it must still resolve material transport, source evolution, and boundary response accurately. The often-used estimate
\begin{equation}
 \frac{\Delta t_{\rm adv}}{\Delta t_{\rm ac}}\sim1+\mathcal M^{-1}
\end{equation}
is local and assumes comparable mesh scales and Courant coefficients. The global material and acoustic minima can occur in different cells. Near zero velocity, the expression is not a finite timestep prescription; temporal error and physical source evolution must set the step.

Preserve the pressure--gravity balance and consistent pressure work at every implicit stage, and compute species transport from the same integrated mass flux as hydrodynamics. Include the open-boundary response in the residual and its derivatives. An acoustic solve communicates across many cells in one step, so freezing ghost values through the nonlinear solve requires justification. Test the hydrostatic base, the heated transition, and the sonic region together; a periodic low-Mach vortex does not test these wind-specific requirements.

Nonlinear convergence error, temporal truncation error, and hydro/source splitting error are distinct. Converging the thermal--composition iteration does not establish the order of the complete hydro/source update. A second-order hydro method combined with an unchanged first-order split can remain first order overall. Refine the complete physical timestep and its source coupling together. Such checks validate a new integrator; they do not add a drift test or a new acceptance condition to the existing campaign's K=3 criterion. A small nonlinear residual must not be used as a substitute for temporal accuracy.

\subsection{Combined cost and the condition for a net gain}
A useful preliminary model separates hydro cost $H$, source cost $Q$, and other cost $O$ in one accepted reference step. Let an implicit step advance $q$ times as much physical time, cost $r_H H$ in hydro, and cost $m_Q Q$ in source evaluations. With unchanged other work in each new step, the idealized speed ratio is
\begin{equation}
 S\simeq\frac{q(H+Q+O)}{r_H H+m_Q Q+O}.
\end{equation}
The gain is positive only if $q(H+Q+O)>r_H H+m_Q Q+O$. Include rejected attempts in the effective costs and account separately for output tied to fixed physical times. This model explains why an acoustic timestep increase may fail to help if each nonlinear residual repeatedly solves radiation and composition. It also explains why source optimization can remain worthwhile when hydro is already cheap. None of $q$, $r_H$, or $m_Q$ has been measured for an EXHALE implicit candidate.

For a local radial stencil with $N$ cells and $b$ coupled variables in each cell, a fixed-width block factorization has cost proportional to $Nb^3$ and storage proportional to $Nb^2$. Thus linear scaling with $N$ alone does not guarantee low cost if all species are included in each block. Compare a hydro block with converged outer source iterations against a more tightly coupled system. Nonlocal column attenuation prevents assuming the exact full Jacobian is block-banded. The decision should follow measured residual and factorization costs, with the same physical equations and error targets in both cases.

\subsection{Concrete order of investigation}
First use existing timers and accepted-step records to measure cost for a common physical interval. Determine whether acoustic-limiting cells are localized and whether chemistry, thermal closure, failed steps, or hydro dominates. The relevant global opportunity is the ratio of \emph{global} material and acoustic timestep limits, not the inverse Mach number of one chosen base cell.

If the acoustic interval dominates, compare a conservative implicit hydro update using a directly factored radial block system with the same full time residual using an acoustic preconditioner. Keep the spherical pressure--gravity balance, mass-consistent species fluxes, and thermal closure in both. Compare backward Euler as a verification baseline and a suitable second-order method for useful accuracy; do not assume the larger stable step wins. For a local hydro stencil with fixed block size, a direct banded factorization scales linearly with radial cell count, so a large multidimensional JFNK framework need not be the starting point.

If source work dominates, first measure repeated rate evaluations, Jacobian construction, factorization, and coupled iterations separately. Test reuse only where the actual equations permit it; consider adaptive subcycling only with a common physical clock and a controlled radiation/composition error. Townsend-style scalar integration is attractive only if a repeatedly used, sufficiently fixed thermal state makes its preparation economical. Preserve nucleus conservation, thermal energy, trace-species resolution, and rejection handling.

For any candidate define $W(\Delta t)$ as measured wall time including rejected attempts. Optimize physical time advanced divided by $W$, subject to profile and temporal accuracy, rather than maximizing $\Delta t$. The Viallet example demonstrates why this distinction matters. These are physical-time integration candidates; obtaining a stationary Newton solution is not required.

\section{How to decide whether it is faster}
Let $q=\Delta t_{\rm new}/\Delta t_{\rm old}$ and let $r$ be the ratio of measured wall cost for one accepted new step to one accepted old step, including failed attempts and amortized setup. Over the same physical interval,
\begin{equation}
 S\simeq q/r.
\end{equation}
For variable steps, measure the wall time needed to reach the same physical end time, including rejected attempts, source iterations, output, and solver setup performed during the run. Exclude compilation, scheduler wait, and unrelated queue delays. Report time to the accepted final K=3 window separately: artificial damping or a changed trajectory can alter when a method appears to settle.

For a conventional optimization, if a fraction $f$ of existing cost is accelerated by a factor $a$, the idealized total gain is $[(1-f)+f/a]^{-1}$. Thus phase fractions are needed before prioritizing source optimization over hydrodynamic redesign. No present bottleneck fraction, allowable implicit step, or speedup is established by this report.

\subsection{Controlled next experiment}
Use identical reviewed initial states, input physics, mesh, boundary conditions, and physical output times. Run the existing explicit hydro integrator at its validated settings as the performance reference, and a temporally refined calculation as the accuracy reference. These serve different purposes: timing only against an unnecessarily small explicit step would overstate the gain. Match CPU model, thread count, compiler optimization, and numerical-library threading; record host load and repeated timing variability. Keep the production campaign unchanged.

First record the limiting cell, local $c_s$, $v$, geometry, accepted steps, rejection reasons, thermal/chemical iterations, phase timings, and physical time advanced. Then evaluate the candidate at several timestep tolerances against a temporally refined explicit calculation over identical physical times. Compare full radial $T$, $\rho$, $v$, and composition profiles, along with mass, nuclei, and energy budgets, the mass flux $4\pi r^2\rho v$, and the sonic-point location when a sonic crossing is present. Include plots and inspect the complete profiles; sample values alone cannot establish agreement.

\noindent\begin{tabularx}{\textwidth}{@{}p{0.29\textwidth}X@{}}
\toprule
Condition to match & Required record for a numerical comparison\\
\midrule
Potential and geometry & Planet/star parameters, tidal prescription, spherical volumes and faces.\\
Domain and boundaries & Both radii, grid, base thermodynamics and composition, outer boundary, initial state.\\
Irradiation & Spectral shape, integrated flux, attenuation, secondary ionization and diffuse contributions.\\
Material physics & Species, reactions, abundances, caloric EOS, cooling/heating and diffusion options.\\
Numerical conditions & Reconstruction, flux formulation, tolerances, thread count, compiler, CPU model, sampling times.\\
\bottomrule
\end{tabularx}

No matched numerical comparison was performed here, so no comparison plots or measured profile agreement are claimed. The table is the required specification for that experiment, not a claim that the supplied papers match EXHALE's inputs.

The campaign's acceptance remains all-cell stability at every stored snapshot over the final K=3 flow-crossing window: $T$ and $\rho$ within 0.004 dex, and $v$ within 1 percent using the prescribed maximum-speed normalization near sign changes and in the quasi-hydrostatic region. Report the measured bands. Do not add a drift-trend gate or wait for a slow chemical process after the specified profiles qualify. Temporal refinement checks numerical fidelity; it does not redefine the campaign criterion or require a stationary certificate.

\section{Recommendation and limits}
Adopt conservative implicit physical-time integration as a research candidate, with Leidi and Viallet providing the most direct additional guidance on hydro integration and solver cost. Dorfi remains a useful early reference. Do not transfer the complete stellar, radiation-diffusion, or scalar-cooling models into EXHALE. First obtain the available phase timings and timestep-limiting diagnostics; then decide whether fixed-grid implicit compressible integration or source-cost reduction has the better measured opportunity. A pressure-only split requires its own derivation and is not selected by these papers alone.

This report establishes applicability conditions, not measured EXHALE performance. No solver source was changed or benchmark run. The principal unresolved quantities are the current bottleneck fractions, global acoustic-to-advective timestep ratio, source-limited accurate step, and nonlinear cost. Until these are measured, a claimed speedup would be speculation.

\begin{thebibliography}{99}
\bibitem{dorfi} E. A. Dorfi (1999), Implicit radiation hydrodynamics for 1D-problems, \emph{Journal of Computational and Applied Mathematics} 109, 153--171. \href{https://doi.org/10.1016/S0377-0427(99)00157-0}{doi:10.1016/S0377-0427(99)00157-0}.
\bibitem{ghosh} K. Ghosh and S. V. G. Menon (2010), Fully implicit 1D radiation hydrodynamics: Validation and verification, \emph{Journal of Computational Physics} 229, 7488--7502. \href{https://doi.org/10.1016/j.jcp.2010.06.031}{doi:10.1016/j.jcp.2010.06.031}.
\bibitem{gardner} D. J. Gardner, J. E. Guerra, F. P. Hamon, D. R. Reynolds, P. A. Ullrich, and C. S. Woodward (2018), Implicit--explicit (IMEX) Runge--Kutta methods for non-hydrostatic atmospheric models, \emph{Geoscientific Model Development} 11, 1497--1515. \href{https://doi.org/10.5194/gmd-11-1497-2018}{doi:10.5194/gmd-11-1497-2018}.
\bibitem{schlegel} M. Schlegel, O. Knoth, M. Arnold, and R. Wolke (2012), Implementation of multirate time integration methods for air pollution modelling, \emph{Geoscientific Model Development} 5, 1395--1405. \href{https://doi.org/10.5194/gmd-5-1395-2012}{doi:10.5194/gmd-5-1395-2012}.

\bibitem{edelmann} P. V. F. Edelmann, L. Horst, J. P. Berberich, R. Andrassy, J. Higl, G. Leidi, C. Klingenberg, and F. K. R\"opke (2021), Well-balanced treatment of gravity in astrophysical fluid dynamics simulations at low Mach numbers, \emph{Astronomy \& Astrophysics} 652, A53. \href{https://doi.org/10.1051/0004-6361/202140653}{doi:10.1051/0004-6361/202140653}.
\bibitem{leidi} G. Leidi, C. Birke, R. Andrassy, J. Higl, P. V. F. Edelmann, G. Wiest, C. Klingenberg, and F. K. R\"opke (2022), A finite-volume scheme for modeling compressible magnetohydrodynamic flows at low Mach numbers in stellar interiors, \emph{Astronomy \& Astrophysics} 668, A143. \href{https://doi.org/10.1051/0004-6361/202244665}{doi:10.1051/0004-6361/202244665}.
\bibitem{bispen} G. Bispen, M. Luk\'a\v{c}ov\'a-Medvid'ov\'a, and L. Yelash (2017), Asymptotic preserving IMEX finite volume schemes for low Mach number Euler equations with gravitation, \emph{Journal of Computational Physics} 335, 222--248. \href{https://doi.org/10.1016/j.jcp.2017.01.020}{doi:10.1016/j.jcp.2017.01.020}.
\bibitem{knoll} D. A. Knoll, W. J. Rider, and G. L. Olson (2001), Nonlinear convergence, accuracy, and time step control in nonequilibrium radiation diffusion, \emph{Journal of Quantitative Spectroscopy and Radiative Transfer} 70, 25--36. \href{https://doi.org/10.1016/S0022-4073(00)00112-6}{doi:10.1016/S0022-4073(00)00112-6}.
\bibitem{viallet13} M. Viallet, I. Baraffe, and R. Walder (2013), Comparison of different nonlinear solvers for 2D time-implicit stellar hydrodynamics, \emph{Astronomy \& Astrophysics} 555, A81. \href{https://doi.org/10.1051/0004-6361/201220725}{doi:10.1051/0004-6361/201220725}.
\bibitem{viallet16} M. Viallet, T. Goffrey, I. Baraffe, D. Folini, C. Geroux, M. V. Popov, J. Pratt, and R. Walder (2016), A Jacobian-free Newton-Krylov method for time-implicit multidimensional hydrodynamics. Physics-based preconditioning for sound waves and thermal diffusion, \emph{Astronomy \& Astrophysics} 586, A153. \href{https://doi.org/10.1051/0004-6361/201527339}{doi:10.1051/0004-6361/201527339}.
\bibitem{townsend} R. H. D. Townsend (2009), An Exact Integration Scheme for Radiative Cooling in Hydrodynamical Simulations, \emph{The Astrophysical Journal Supplement Series} 181, 391--397. \href{https://doi.org/10.1088/0067-0049/181/2/391}{doi:10.1088/0067-0049/181/2/391}.
\bibitem{esenturk} E. Esent\"urk, N. L. Abraham, S. Archer-Nicholls, C. Mitsakou, P. Griffiths, A. Archibald, and J. Pyle (2018), Quasi-Newton methods for atmospheric chemistry simulations: implementation in UKCA UM vn10.8, \emph{Geoscientific Model Development} 11, 3089--3108. \href{https://doi.org/10.5194/gmd-11-3089-2018}{doi:10.5194/gmd-11-3089-2018}.
\bibitem{cariolle} D. Cariolle, P. Moinat, H. Teyss\`edre, L. Giraud, B. Josse, and F. Lef\`evre (2017), ASIS v1.0: an adaptive solver for the simulation of atmospheric chemistry, \emph{Geoscientific Model Development} 10, 1467--1485. \href{https://doi.org/10.5194/gmd-10-1467-2017}{doi:10.5194/gmd-10-1467-2017}.
\end{thebibliography}
\end{document}
